\section{Related Work}

\subsection{ZK Proof Systems for Blockchain State}

Groth16~\cite{groth2016size} remains the most widely deployed SNARK
construction for blockchain applications due to its constant-size proofs
(3 group elements, approximately 128~bytes uncompressed) and fast
verification. Its main limitation---a per-circuit trusted setup---is
acceptable in enterprise deployments where the operator controls the
proving infrastructure. PLONK~\cite{gabizon2019plonk} offers a universal
setup but produces larger proofs and slower verification. For the
validium use case, where proof generation latency matters more than
setup flexibility, Groth16 is the established choice.

The Poseidon hash function~\cite{grassi2021poseidon} was designed
specifically for arithmetic circuits over prime fields, achieving
approximately 240~R1CS constraints per 2-to-1 hash in Circom~\cite{elamrani2020merkle}---orders of magnitude fewer than SHA-256
(26,170 constraints) or Keccak-256 (144,830 constraints)~\cite{bkomuves2023hash}.
Poseidon2~\cite{grassi2023poseidon2} further reduces constraint counts
but has not yet been integrated into the standard Circom library
(circomlib). Recent cryptanalysis~\cite{grassi2025grobner} and the
Ethereum Foundation's ongoing Poseidon Cryptanalysis
Initiative~\cite{ethereum2024poseidon} have confirmed adequate security
margins for the parameter sets used in practice.

\subsection{Sparse Merkle Trees in ZK Systems}

Sparse Merkle Trees~\cite{baylina2019smt} provide deterministic,
collision-free key-to-leaf mappings suitable for authenticated state
management. The Iden3/circomlib SMT implementation~\cite{iden3smt}
offers \texttt{SMTVerifier} and \texttt{SMTProcessor} circuit templates,
with the latter handling insert, update, and delete operations at
approximately $1{,}000 + 250d$ constraints for a tree of depth~$d$~\cite{elamrani2020merkle}.
Our circuit takes a different approach: rather than using \texttt{SMTProcessor},
we implement explicit dual-path Merkle verification (old root + new root)
to keep the circuit minimal and update-only.

Binary Merkle path verification in circomlib requires 219~constraints per
tree level, combining a Poseidon(2) hash with a Mux1
selector~\cite{elamrani2020merkle}. Quinary Poseidon trees reduce the
number of levels (capacity $5^d$ vs $2^d$) but increase per-level cost to
approximately 441~constraints, yielding no net benefit for fixed-capacity
requirements.

\subsection{Production ZK State Transition Systems}

Table~\ref{tab:production-comparison} summarizes the constraint characteristics
of production systems that incorporate Merkle-based state transitions.

\begin{table}[htbp]
\centering
\caption{Constraint comparison with production ZK systems}
\label{tab:production-comparison}
\begin{tabular}{@{}lrrr@{}}
\toprule
\textbf{System} & \textbf{Per-Tx} & \textbf{Batch} & \textbf{Notes} \\
\midrule
Tornado Cash~\cite{tornadocash2019} & 5--7K & 1 & Merkle + nullifier \\
Semaphore v4~\cite{semaphore2024v4} & $\sim$5K & 1 & Identity proof \\
Hermez rollup~\cite{hermez2021circuits} & $\sim$58K & 2048 & Full tx logic \\
Polygon zkEVM~\cite{polygon2024zkevm} & --- & var. & STARK+Groth16 \\
\textbf{This work} & \textbf{34.3K} & \textbf{4--16} & \textbf{SMT update only} \\
\bottomrule
\end{tabular}
\par\vspace{1mm}
\parbox{\columnwidth}{\footnotesize Per-tx constraints at depth 32. Hermez includes EdDSA signature
verification, token transfer, and fee processing. This work measures only
Merkle-based state transition cost.}
\end{table}
\FloatBarrier

Hermez~\cite{hermez2021circuits} uses circomlib's \texttt{SMTProcessor}
within a much larger circuit that includes EdDSA signature verification,
token balance updates, and fee distribution, resulting in approximately
58,000~constraints per transaction. Their 2048-transaction circuit
exceeds 118~million constraints and requires a custom C++ prover.
Our per-transaction cost of 34,254~constraints at depth~32 is 0.59$\times$
the Hermez figure, consistent with the fact that our circuit omits
signature and token logic.

Tornado Cash~\cite{tornadocash2019} implements a single Merkle inclusion
proof at depth~20 with a nullifier check, requiring approximately
5,000--7,000~constraints. Semaphore~v4~\cite{semaphore2024v4} achieves
similar constraint counts for identity set membership proofs. Both systems
prove membership in a single tree without state modification, making
direct comparison with our update circuit inappropriate; however, they
establish the baseline cost of a single Merkle path verification.

\subsection{Groth16 Proving Performance}

Proving time is ultimately the binding constraint for enterprise batch
processing. Table~\ref{tab:prover-benchmarks} summarizes published prover
benchmarks.

\begin{table}[htbp]
\centering
\caption{Published Groth16 proving time benchmarks}
\label{tab:prover-benchmarks}
\begin{tabular}{@{}lrrl@{}}
\toprule
\textbf{Prover} & \textbf{Constraints} & \textbf{Time (s)} & \textbf{Hardware} \\
\midrule
Rapidsnark~\cite{rapidsnark2024} & 157K & 1.0 & 32-core server \\
Rapidsnark~\cite{rapidsnark2024} & 1.2M & 3.0 & 32-core server \\
Tachyon~\cite{tachyon2024} & 157K & 0.4 & 32-core server \\
Tachyon~\cite{tachyon2024} & 1.2M & 2.55 & 32-core server \\
SnarkJS~\cite{snarkjs} & $\sim$15K & $\sim$15 & Commodity \\
SnarkJS~\cite{snarkjs} & $\sim$100K & $\sim$60 & Commodity \\
\bottomrule
\end{tabular}
\par\vspace{1mm}
\parbox{\columnwidth}{\footnotesize Rapidsnark and Tachyon benchmarks from \cite{tachyon2024}.
SnarkJS estimates from zk-email and RSA circuit reports~\cite{zkmopro2024}.}
\end{table}
\FloatBarrier

Rapidsnark achieves approximately 2.5~$\mu$s per constraint at 1.2~million
constraints on a 32-core AMD Threadripper. SnarkJS, being a pure JavaScript
implementation, is 4--10$\times$ slower. GPU-accelerated provers such as
ICICLE-Snark~\cite{ingonyama2024icicle} report further 3--5$\times$
speedups. These benchmarks establish the range of proving times that our
constraint counts imply.
